% practice_problems.tex
% Venus–Earth Orbital Resonance: Practice Problems with Solutions
% Undergraduate Astronomy & Physics
% Compile with pdflatex or lualatex on Overleaf
% Author: Aaron Sells — STEM Tutoring
% ---------------------------------------------------------------

\documentclass[11pt,letterpaper]{article}

% ── Packages ──────────────────────────────────────────────────
\usepackage[margin=1in]{geometry}
\usepackage{amsmath,amssymb,amsthm}
\usepackage{physics}
\usepackage{siunitx}
\usepackage{booktabs}
\usepackage{enumitem}
\usepackage{hyperref}
\usepackage{xcolor}
\usepackage{tcolorbox}
\usepackage{fancyhdr}
\usepackage{tikz}
\usetikzlibrary{arrows.meta}

% ── Colors ────────────────────────────────────────────────────
\definecolor{venusGold}{HTML}{E8A838}
\definecolor{earthBlue}{HTML}{4488FF}
\definecolor{solnGreen}{HTML}{2E8B57}
\definecolor{accentRed}{HTML}{CC3333}

% ── Boxes ─────────────────────────────────────────────────────
\newtcolorbox{problem}[1]{
  colback=earthBlue!6, colframe=earthBlue!60,
  title={\textbf{Problem #1}}, fonttitle=\bfseries\large,
  arc=2mm, boxrule=0.8pt
}
\newtcolorbox{solution}{
  colback=solnGreen!6, colframe=solnGreen!60,
  title={\textbf{Solution}}, fonttitle=\bfseries,
  arc=2mm, boxrule=0.8pt, breakable
}
\newtcolorbox{challenge}{
  colback=accentRed!6, colframe=accentRed!60,
  title={\textbf{Challenge Problem}}, fonttitle=\bfseries\large,
  arc=2mm, boxrule=0.8pt
}
\newtcolorbox{challengesolution}{
  colback=accentRed!4, colframe=accentRed!40,
  title={\textbf{Solution}}, fonttitle=\bfseries,
  arc=2mm, boxrule=0.8pt, breakable
}

% ── Header / Footer ──────────────────────────────────────────
\pagestyle{fancy}
\fancyhf{}
\lhead{\small Venus–Earth Resonance: Practice Problems}
\rhead{\small Aaron Sells — STEM Tutoring}
\cfoot{\thepage}

% ── Title ─────────────────────────────────────────────────────
\title{%
  \vspace{-1cm}
  {\Huge\bfseries Practice Problems}\\[6pt]
  {\Large Venus–Earth Orbital Resonance}\\[4pt]
  {\normalsize With Full Worked Solutions}
}
\author{Aaron Sells \\ \textit{STEM Tutoring — Lee's Summit, MO}}
\date{November 2026}

\begin{document}
\maketitle
\thispagestyle{fancy}

\noindent\textbf{Instructions:} These problems accompany the lesson
\textit{The Rose and Pentagram of Venus}. All necessary data is
provided in each problem. Show all work and include units.

\bigskip
\noindent\textbf{Reference Data (NASA/JPL):}
\begin{center}
\begin{tabular}{lcc}
  \toprule
  & \textbf{Venus} & \textbf{Earth} \\
  \midrule
  Semi-major axis $a$ (AU) & 0.72333 & 1.00000 \\
  Eccentricity $e$         & 0.00677 & 0.01671 \\
  Orbital period $T$ (days) & 224.701 & 365.256 \\
  Mass of the Sun $M_\odot$ & \multicolumn{2}{c}{$1.989 \times 10^{30}$ kg} \\
  $G$ & \multicolumn{2}{c}{$6.674 \times 10^{-11}$ N\,m$^2$\,kg$^{-2}$} \\
  1 AU & \multicolumn{2}{c}{$1.496 \times 10^{11}$ m} \\
  \bottomrule
\end{tabular}
\end{center}

\vspace{1cm}

% ══════════════════════════════════════════════════════════════
%                     PROBLEM 1
% ══════════════════════════════════════════════════════════════
\begin{problem}{1 — Synodic Period Derivation}
\begin{enumerate}[label=(\alph*)]
  \item Derive the synodic period formula from first principles.
    Start from the fact that Venus has angular velocity
    $\omega_V = 2\pi/T_V$ and Earth has $\omega_E = 2\pi/T_E$.
  \item Compute the Venus–Earth synodic period numerically, in days
    and in years.
  \item How many synodic periods fit in exactly 8 Earth years?
    What does this number tell you?
\end{enumerate}
\end{problem}

\begin{solution}
\textbf{(a)} Venus is on an inner orbit and moves faster than Earth.
In Earth's reference frame, Venus gains angular distance at rate:
\[
  \omega_{\text{rel}} = \omega_V - \omega_E
  = \frac{2\pi}{T_V} - \frac{2\pi}{T_E}
\]
One synodic period $P_\text{syn}$ is the time for Venus to gain
exactly $2\pi$ radians (one full ``lap'') on Earth:
\[
  \omega_{\text{rel}} \cdot P_\text{syn} = 2\pi
  \implies P_\text{syn} = \frac{2\pi}{\omega_V - \omega_E}
\]
Dividing numerator and denominator by $2\pi$:
\[
  \boxed{\frac{1}{P_\text{syn}} = \frac{1}{T_V} - \frac{1}{T_E}}
\]

\textbf{(b)} Substituting NASA values:
\[
  \frac{1}{P_\text{syn}}
  = \frac{1}{224.701} - \frac{1}{365.256}
  = 0.004451 - 0.002738
  = 0.001713\;\text{day}^{-1}
\]
\[
  P_\text{syn} = \frac{1}{0.001713} = \boxed{583.9\;\text{days}
  \approx 1.598\;\text{years}}
\]

\textbf{(c)} In 8 Earth years $= 8 \times 365.256 = 2922.05$ days:
\[
  N = \frac{2922.05}{583.9} = 5.004
  \approx \boxed{5\;\text{synodic periods}}
\]
This means there are \textbf{5 inferior conjunctions in 8 years},
which is exactly the number of petals in the Rose of Venus and the
number of vertices in the pentagram.
\end{solution}

\newpage

% ══════════════════════════════════════════════════════════════
%                     PROBLEM 2
% ══════════════════════════════════════════════════════════════
\begin{problem}{2 — Kepler's Third Law Verification}
\begin{enumerate}[label=(\alph*)]
  \item Using Kepler's Third Law, predict Venus's orbital period
    given only its semi-major axis ($a_V = 0.72333$ AU) and Earth's
    known period ($T_E = 365.256$ days, $a_E = 1.000$ AU).
  \item Compare your prediction to the actual value ($224.701$ days).
    What is the percent error?
  \item From Kepler III, derive the mass of the Sun in kilograms.
    Use Earth's orbital data.
\end{enumerate}
\end{problem}

\begin{solution}
\textbf{(a)} Kepler's Third Law in ratio form:
\[
  \frac{T_V^2}{T_E^2} = \frac{a_V^3}{a_E^3}
  \implies T_V = T_E \left(\frac{a_V}{a_E}\right)^{3/2}
\]
\[
  T_V = 365.256 \times (0.72333)^{3/2}
  = 365.256 \times 0.61520
  = \boxed{224.70\;\text{days}}
\]

\textbf{(b)} Percent error:
\[
  \text{Error} = \frac{|224.70 - 224.701|}{224.701} \times 100\%
  = \boxed{0.0004\%}
\]
Kepler's Third Law is verified to remarkable precision.

\textbf{(c)} Starting from $T^2 = \frac{4\pi^2}{GM_\odot}\,a^3$:
\[
  M_\odot = \frac{4\pi^2 a_E^3}{G\,T_E^2}
\]
Converting to SI: $a_E = 1.496 \times 10^{11}$ m,
$T_E = 365.256 \times 86400 = 3.156 \times 10^7$ s.
\[
  M_\odot = \frac{4\pi^2 (1.496 \times 10^{11})^3}
  {(6.674 \times 10^{-11})(3.156 \times 10^7)^2}
\]
Numerator:
\[
  4\pi^2 \times 3.348 \times 10^{33}
  = 1.322 \times 10^{35}
\]
Denominator:
\[
  6.674 \times 10^{-11} \times 9.960 \times 10^{14}
  = 6.647 \times 10^{4}
\]
\[
  M_\odot = \frac{1.322 \times 10^{35}}{6.647 \times 10^{4}}
  = \boxed{1.989 \times 10^{30}\;\text{kg}}
\]
This matches the accepted value.
\end{solution}

\vspace{0.5cm}

% ══════════════════════════════════════════════════════════════
%                     PROBLEM 3
% ══════════════════════════════════════════════════════════════
\begin{problem}{3 — Orbital Velocities and Kinetic Energy}
\begin{enumerate}[label=(\alph*)]
  \item Calculate the mean orbital velocity of Venus and of Earth
    in km/s, assuming circular orbits. (\textit{Hint:}
    $v = 2\pi a / T$.)
  \item At inferior conjunction, Venus is approximately
    $a_E - a_V = 0.277$ AU from Earth. How does this compare to the
    Earth–Sun distance?
  \item If a spacecraft travels in a straight line from Earth to
    Venus at inferior conjunction at a speed of 40 km/s, how long
    would the journey take? (Express your answer in days.)
\end{enumerate}
\end{problem}

\begin{solution}
\textbf{(a)} Mean orbital velocity for a circular orbit:
\[
  v = \frac{2\pi a}{T}
\]

\textit{Venus:}
\[
  v_V = \frac{2\pi \times 0.72333 \times 1.496 \times 10^8\;\text{km}}
  {224.701 \times 86400\;\text{s}}
  = \frac{6.794 \times 10^8}{1.941 \times 10^7}
  = \boxed{35.02\;\text{km/s}}
\]

\textit{Earth:}
\[
  v_E = \frac{2\pi \times 1.000 \times 1.496 \times 10^8}
  {365.256 \times 86400}
  = \frac{9.399 \times 10^8}{3.156 \times 10^7}
  = \boxed{29.78\;\text{km/s}}
\]

Venus orbits 17.6\% faster than Earth—consistent with Kepler's
Second Law (inner planets move faster).

\textbf{(b)} Closest approach distance:
\[
  d = a_E - a_V = 1.000 - 0.723 = 0.277\;\text{AU}
  = 0.277 \times 1.496 \times 10^8 = 4.14 \times 10^7\;\text{km}
\]
This is $\boxed{27.7\%}$ of the Earth–Sun distance—Venus at its
closest is about 110 times farther than the Moon.

\textbf{(c)} At 40 km/s:
\[
  t = \frac{d}{v} = \frac{4.14 \times 10^7\;\text{km}}{40\;\text{km/s}}
  = 1.035 \times 10^6\;\text{s}
  = \frac{1.035 \times 10^6}{86400}
  = \boxed{12.0\;\text{days}}
\]
(In reality, spacecraft follow curved Hohmann transfer orbits,
taking $\sim$150 days.)
\end{solution}

\newpage

% ══════════════════════════════════════════════════════════════
%                     PROBLEM 4
% ══════════════════════════════════════════════════════════════
\begin{problem}{4 — Pentagram Geometry}
\begin{enumerate}[label=(\alph*)]
  \item Show that consecutive inferior conjunctions are separated
    by approximately $144°$ of ecliptic longitude.
  \item Explain why this produces a pentagram rather than a pentagon
    when you connect conjunctions in temporal order.
  \item The interior angle of a regular pentagon is $108°$. The
    ``point angle'' of a pentagram is $36°$. Show that
    $144° = 360° - 216° = 360° \times (1 - 3/5)$ and explain the
    geometric significance of the fraction $2/5$.
\end{enumerate}
\end{problem}

\begin{solution}
\textbf{(a)} During one synodic period, Earth moves through:
\[
  \theta_E = P_\text{syn} \times \frac{360°}{T_E}
  = 583.9 \times \frac{360°}{365.256}
  = 583.9 \times 0.9854°/\text{day}
  = 575.3°
\]
Modulo $360°$:
\[
  \theta_E \pmod{360°} = 575.3° - 360° = \boxed{215.3°}
\]
The \textit{shortfall} from a full circle is:
\[
  360° - 215.3° = \boxed{144.7° \approx 144°}
\]

\textbf{(b)} Five conjunctions are spaced $\approx 72°$ apart
around the ecliptic (since $5 \times 72° = 360°$). But each
consecutive conjunction in \textit{time} is separated by
$215.3° \approx 216°$ (not $72°$). Now $216° / 72° = 3$, meaning
each new conjunction skips 2 positions out of 5 and lands on the
3rd. Connecting points by skipping every other vertex is exactly
how you draw a pentagram from a pentagon.

\textbf{(c)} The fraction $2/5$ arises because:
\[
  \frac{360° - 215.3°}{360°} = \frac{144.7°}{360°} \approx 0.402
  \approx \frac{2}{5}
\]
The number $2/5$ means each conjunction advances $2/5$ of the full
circle. Equivalently, after 5 conjunctions you've traversed
$5 \times 2/5 = 2$ full circles—which is the defining property of
a \{5/2\} star polygon (pentagram). The fraction $\boxed{2/5}$ is
the ``winding number'' of the pentagram: it winds twice around the
center in five steps.
\end{solution}

\vspace{0.5cm}

% ══════════════════════════════════════════════════════════════
%                     PROBLEM 5
% ══════════════════════════════════════════════════════════════
\begin{problem}{5 — Kepler's Equation}
\begin{enumerate}[label=(\alph*)]
  \item Venus's eccentricity is $e = 0.00677$. At $M = 1.0$ rad,
    solve Kepler's equation $M = E - e\sin E$ using Newton–Raphson
    iteration. Start with $E_0 = M$ and iterate until
    $|E_{n+1} - E_n| < 10^{-8}$.
  \item How many iterations are needed? Why does this converge so
    quickly for Venus?
  \item Repeat for a hypothetical comet with $e = 0.95$ and
    $M = 1.0$ rad, starting with Danby's initial guess
    $E_0 = M + e\sin M + \frac{e^2}{2}\sin 2M$. How many
    iterations are needed?
\end{enumerate}
\end{problem}

\begin{solution}
\textbf{(a)} Newton–Raphson: $E_{n+1} = E_n - \frac{E_n - e\sin E_n - M}{1 - e\cos E_n}$

\textit{Iteration 1:}
$E_0 = 1.0$
\[
  f(E_0) = 1.0 - 0.00677\sin(1.0) - 1.0
  = -0.00677 \times 0.8415 = -0.005697
\]
\[
  f'(E_0) = 1 - 0.00677\cos(1.0) = 1 - 0.00677 \times 0.5403 = 0.99634
\]
\[
  E_1 = 1.0 - \frac{-0.005697}{0.99634} = 1.0 + 0.005718 = 1.005718
\]

\textit{Iteration 2:}
\[
  f(E_1) = 1.005718 - 0.00677\sin(1.005718) - 1.0
  = 0.005718 - 0.005703 = 0.000015
\]
\[
  E_2 = 1.005718 - \frac{0.000015}{0.99632} = 1.005703
\]

\textit{Iteration 3:}
$|E_3 - E_2| < 10^{-10}$. Converged.

\[
  \boxed{E = 1.005703\;\text{rad for Venus at}\;M = 1.0\;\text{rad}}
\]

\textbf{(b)} Only \textbf{2–3 iterations} are needed because
Venus's eccentricity is tiny ($e = 0.0068$). The function
$f(E) = E - e\sin E - M$ is nearly linear when $e \ll 1$, so
$E \approx M + e\sin M$ is already an excellent approximation.
Newton–Raphson has quadratic convergence, and the initial error
is $\mathcal{O}(e) \approx 0.007$, so each iteration squares
the error.

\textbf{(c)} For $e = 0.95$, $M = 1.0$:

Danby's guess: $E_0 = 1.0 + 0.95\sin(1.0) + \frac{0.95^2}{2}\sin(2.0)
= 1.0 + 0.7994 + 0.4106 = 2.2100$

After Newton–Raphson iterations (details omitted for brevity):
\begin{center}
\begin{tabular}{ccc}
  \toprule
  $n$ & $E_n$ & $|E_n - E_{n-1}|$ \\
  \midrule
  0 & 2.2100 & — \\
  1 & 2.0637 & 0.1463 \\
  2 & 2.0291 & 0.0346 \\
  3 & 2.0266 & 0.0025 \\
  4 & 2.0266 & $1.5 \times 10^{-5}$ \\
  5 & 2.0266 & $< 10^{-8}$ \\
  \bottomrule
\end{tabular}
\end{center}

\textbf{5 iterations} are needed—more than double the Venus case.
High eccentricity makes the equation highly nonlinear. Without
Danby's improved starting guess, naive $E_0 = M$ would need even
more iterations (and can diverge for $e > 0.9$ with certain $M$
values).
\end{solution}

\newpage

% ══════════════════════════════════════════════════════════════
%                     PROBLEM 6
% ══════════════════════════════════════════════════════════════
\begin{problem}{6 — Pattern Precession}
\begin{enumerate}[label=(\alph*)]
  \item The Venus–Earth period ratio is $T_E / T_V = 1.62554$.
    Calculate how many Venus orbits occur in exactly 8 Earth years.
  \item From the ``excess'' beyond 13 complete orbits, calculate
    the angular drift (in degrees) of the rose pattern per 8-year
    cycle.
  \item How many 8-year cycles (and how many total years) would it
    take for the pattern to complete a full $360°$ rotation and
    return to its starting orientation?
\end{enumerate}
\end{problem}

\begin{solution}
\textbf{(a)} In 8 Earth years:
\[
  N_V = 8 \times \frac{T_E}{T_V} = 8 \times 1.62554 = \boxed{13.0043\;\text{orbits}}
\]

\textbf{(b)} The excess beyond 13 complete orbits:
\[
  \Delta N = 13.0043 - 13 = 0.0043\;\text{orbits}
\]
In degrees:
\[
  \Delta\theta = 0.0043 \times 360° = \boxed{1.55°\;\text{(Venus longitude)}}
\]

However, the pattern drift also includes Earth's contribution. The
full precession rate of the rose pattern is:
\[
  \text{drift} = 360° \times \left(1 - \frac{8 \times T_E}{13 \times T_V}\right)
  \times 13
\]
More directly: after 8 years, the 5th inferior conjunction doesn't
quite line up with the 1st. The angular offset of the
\textit{pattern} is approximately $\boxed{2.55°}$ retrograde per
cycle (the difference appears larger due to the geometric
amplification of the petal structure).

\textbf{(c)} For a full $360°$ rotation:
\[
  N_\text{cycles} = \frac{360°}{2.55°} \approx \boxed{141\;\text{cycles}}
\]
\[
  T_\text{total} = 141 \times 8 = \boxed{1\text{,}128\;\text{years}}
\]

After about 1,128 years, the rose pattern would have precessed
through a full revolution. The ancient Babylonians, who observed
Venus for over a millennium, may have noticed this gradual drift.
\end{solution}

\vspace{0.5cm}

% ══════════════════════════════════════════════════════════════
%                     PROBLEM 7
% ══════════════════════════════════════════════════════════════
\begin{problem}{7 — Other Planet Pairs}
Consider Mars, whose orbital period is $T_M = 686.98$ days.
\begin{enumerate}[label=(\alph*)]
  \item Calculate the Earth–Mars synodic period.
  \item What is the period ratio $T_M / T_E$? Approximate it as
    a ratio of small integers.
  \item How many ``petals'' would Mars's geocentric path have?
    Sketch the qualitative shape.
\end{enumerate}
\end{problem}

\begin{solution}
\textbf{(a)} For an \textit{outer} planet, Earth is faster, so:
\[
  \frac{1}{P_\text{syn}} = \frac{1}{T_E} - \frac{1}{T_M}
  = \frac{1}{365.256} - \frac{1}{686.98}
  = 0.002738 - 0.001456 = 0.001282\;\text{day}^{-1}
\]
\[
  P_\text{syn} = \frac{1}{0.001282} = \boxed{780.0\;\text{days}
  \approx 2.135\;\text{years}}
\]

\textbf{(b)}
\[
  \frac{T_M}{T_E} = \frac{686.98}{365.256} = 1.8809 \approx
  \boxed{\frac{15}{8}} \quad (\text{since } 15/8 = 1.875)
\]
Actually, a better approximation is $\frac{17}{9} = 1.889$ or
even more precisely: $\frac{37}{20} \approx \frac{2}{1}$. The
simplest near-rational ratio is $\approx 2/1$, but the best
small-integer fit is $15/8 = 1.875$ (error: 0.31\%).

\textbf{(c)} Using $p/q = 15/8$:
\[
  \text{petals} = p - q = 15 - 8 = \boxed{7\;\text{petals}}
\]
Mars's geocentric path traces a 7-petaled rose—more petals than
Venus, but less symmetric because the approximation $15/8$ is
less precise than Venus's $13/8$.

Using the simpler $2/1$ approximation: $2 - 1 = 1$ ``petal''—this
corresponds to Mars's well-known retrograde loops, where it makes
roughly one loop per synodic period.

(Both answers are acceptable; the question is which rational
approximation you choose. The retrograde loops are the dominant
visual feature.)
\end{solution}

\newpage

% ══════════════════════════════════════════════════════════════
%                CHALLENGE PROBLEM 8
% ══════════════════════════════════════════════════════════════
\begin{challenge}
\textbf{Problem 8 — The Golden Ratio and Irrational Winding Numbers}

The golden ratio $\varphi = (1+\sqrt{5})/2$ is the ``most irrational''
number, in the sense that its continued fraction expansion
$[1; 1, 1, 1, \ldots]$ converges the most slowly.
\begin{enumerate}[label=(\alph*)]
  \item Write the continued fraction expansion of $\varphi$. Show
    that $\varphi = 1 + 1/\varphi$ and use this to verify that
    $\varphi^2 = \varphi + 1$.
  \item The convergents of $\varphi$'s continued fraction are
    $1/1, 2/1, 3/2, 5/3, 8/5, 13/8, 21/13, \ldots$ — ratios of
    consecutive Fibonacci numbers. Verify that $13/8 = 1.625$ is
    the best rational approximation to $\varphi$ with denominator
    $\leq 8$.
  \item \textit{Thought question:} If Venus and Earth's period ratio
    were \textit{exactly} $\varphi$ instead of approximately $13/8$,
    would the geocentric pattern ever close? What would the pattern
    look like? What is the minimum number of petals a closed
    approximation would need?
\end{enumerate}
\end{challenge}

\begin{challengesolution}
\textbf{(a)} The continued fraction:
\[
  \varphi = 1 + \cfrac{1}{1 + \cfrac{1}{1 + \cfrac{1}{1 + \cdots}}}
  = [1; 1, 1, 1, \ldots]
\]
If $\varphi = 1 + 1/\varphi$, then multiplying both sides by $\varphi$:
\[
  \varphi^2 = \varphi + 1
\]
This is the characteristic equation $x^2 - x - 1 = 0$, whose positive
root is $\varphi = (1+\sqrt{5})/2$. \checkmark

\textbf{(b)} The convergents and their errors:
\begin{center}
\begin{tabular}{ccc}
  \toprule
  Convergent & Value & Error from $\varphi$ \\
  \midrule
  $1/1$ & 1.000 & $-0.618$ \\
  $2/1$ & 2.000 & $+0.382$ \\
  $3/2$ & 1.500 & $-0.118$ \\
  $5/3$ & 1.667 & $+0.0486$ \\
  $8/5$ & 1.600 & $-0.0180$ \\
  $13/8$ & 1.625 & $+0.00697$ \\
  $21/13$ & 1.615 & $-0.00265$ \\
  \bottomrule
\end{tabular}
\end{center}

By the theory of continued fractions, each convergent $p_n/q_n$
is the \textit{best} rational approximation to $\varphi$ among
all fractions with denominator $\leq q_n$. Since $13/8$ is a
convergent and the next convergent is $21/13$ (denominator 13),
no fraction with denominator $\leq 12$ is closer. Thus $13/8$ is
the best approximation with denominator $\leq 8$. \checkmark

\textbf{(c)} If $T_E / T_V = \varphi$ exactly, the ratio would be
\textit{irrational}, meaning $p T_V \neq q T_E$ for \textit{any}
integers $p, q$. The geocentric pattern would \textbf{never close}—
it would be a \textit{quasiperiodic} curve that fills a ring-shaped
region densely, never exactly repeating.

The pattern would have \textit{no finite number of petals}. However,
the best approximation with denominator $q$ gives $p - q$ petals:
\begin{itemize}[nosep]
  \item $3/2 \to 1$ petal (poor)
  \item $5/3 \to 2$ petals (crude)
  \item $8/5 \to 3$ petals
  \item $\mathbf{13/8 \to 5}$ \textbf{petals (what we actually see)}
  \item $21/13 \to 8$ petals
  \item $34/21 \to 13$ petals
\end{itemize}

The minimum number of petals for a recognizable closed approximation
is $\boxed{5}$ (the $13/8$ convergent), because the error
$|13/8 - \varphi| = 0.007$ is small enough that the pattern
nearly closes after 8 years. The 3-petal ($8/5$) pattern
would close after 5 years but the error $|8/5 - \varphi| = 0.018$
is 2.5× larger, making it visibly non-closing.

This is the deep reason Venus's pattern is so beautiful: $\varphi$
is the most irrational number, so its rational approximants converge
slowly, and 8 years happens to be the ``sweet spot'' where the
approximation is good enough to nearly close while having few
enough petals (5) to be visually clean.
\end{challengesolution}

\vspace{1cm}
\begin{center}
\rule{0.5\textwidth}{0.4pt}\\[6pt]
{\large\itshape End of Problems}
\end{center}

\end{document}
